\section*{Introduction}
\addcontentsline{toc}{section}{Introduction}

Let $A$ be a central simple algebra over a field $F$.
In degrees $i=1,2$, the groups $K_i(A)$ are compared with $K_i(F)$ by the reduced norms $\Nrd_i\colon K_i(A)\to K_i(F)$ \cite{MerkurjevSuslin1982,Suslin1987}.
We study these maps over nonarchimedean local fields, global fields, and two- and three-dimensional local fields.


\subsection*{Background}

Let $F$ be a local or global field.
Merkurjev and Suslin conjectured that the sequence
\begin{equation}\label{eq:intro-MS}
 0\longrightarrow K_2(A)\xrightarrow{\ \Nrd_2\ }K_2(F)
 \xrightarrow{\ (s_v)_{v}\ }(\Z/2)^{S_{\mathbb R}}\longrightarrow0
\end{equation}
is exact, where $S_{\mathbb R}$ is the set of real places of $F$ at which $A$ is ramified and $s_v$ is the Hilbert symbol at $v$ \cite[Remark 17.5]{MerkurjevSuslin1982}.
They proved~\eqref{eq:intro-MS} when the index of $A$ is squarefree and prime to the characteristic of $F$ \cite[Theorem 17.4]{MerkurjevSuslin1982}.
Merkurjev and Suslin and, independently, Kahn and Levine proved the injectivity of $\Nrd_2$ for $A$ of squarefree index over every field of characteristic zero \cite[Theorem 2.3]{MerkurjevSuslin2010}, \cite[Corollary 2]{KahnLevine2010}.
For global fields, Bak and Rehmann identified the image of $\Nrd_2$ with the kernel of $(s_v)_v$ \cite[\S2, Theorem 3]{BakRehmann1984}, so the main point of~\eqref{eq:intro-MS} is the injectivity of $\Nrd_2$.
This injectivity over global fields has sometimes been attributed to \cite{BakRehmann1984}; as pointed out in \cite[Remark 10.4(3)]{BHRAdeles}, however, \cite{BakRehmann1984} does not claim it.

The corresponding statements in degree one are classical \cite{NakayamaMatsushima1943,Wang1950}, \cite[Corollary 4.10(iii)]{HazratWadsworth}, \cite[Theorem 33.15]{ReinerMaximalOrders}, \cite[Corollary 2]{Yanchevskii1975}.


Let $F$ be a nonarchimedean local field with residue characteristic $p$, and let $D$ be the division algebra in the Brauer class of $A$.
Suslin and Yufryakov showed that, if $\deg D$ is prime to $p$, then for every prime $\ell\ne p$ and every $j\ge1$ there is an isomorphism $\pi_jK(D)^\wedge_\ell\cong\pi_jK(F)^\wedge_\ell$ whose product with $\deg D$ is induced by the norm $N_{D/F}$ \cite[Theorem 3]{SuslinYufryakov}, and Hesselholt, Larsen and Lindenstrauss proved the analogous statement for $\ell=p$ when $p$ is odd \cite[Theorem A]{HesselholtLarsenLindenstrauss2020}.

Clausen described the Artin map of class field theory through the $K$-theory of locally compact $F$-modules \cite{ClausenArtin}.
Braunling, Henrard and van Roosmalen developed a noncommutative analogue, which relates the id\`eles of $A$ to $K_1$ of the category $\operatorname{LCA}_A$ of locally compact $A$-modules \cite{BHRAdeles}.
Reduced norms are closely related to noncommutative $K$-theoretic class field theory; for example, the Merkurjev--Suslin conjecture implies the isomorphism $K_2(\operatorname{LCA}_A)\cong K_2(\operatorname{LCA}_F)$ of \cite[Conjecture 10.2]{BHRAdeles}; see \cite[Theorem 10.5]{BHRAdeles}.


\subsection*{Main results}

\begin{introthm}[Merkurjev--Suslin conjecture, Theorems~\ref{thm:paper-local-norm} and~\ref{thm:paper-global-norm}]\label{thm:intro-K2}
Let $F$ be a nonarchimedean local field or a global field and $A$ a central simple algebra over $F$.
Then $\Nrd_2$ is injective.
In particular, the sequence
\[
 0\longrightarrow K_2(A)\xrightarrow{\ \Nrd_2\ }K_2(F)
 \xrightarrow{\ (s_v)_{v}\ }(\Z/2)^{S_{\mathbb R}}\longrightarrow0
\]
is exact.
\end{introthm}


Hesselholt, Larsen and Lindenstrauss remark that it is not known whether their reduced norm isomorphism is induced by a map of $K$-theory spectra, and that such a map may well not exist \cite[Introduction]{HesselholtLarsenLindenstrauss2020}; compare \cite[Proposition 4]{Merkurjev1995}.
We construct a map of truncated $K$-theory spectra that induces the reduced norms in degrees $1$ and $2$.
Moreover, over local fields we obtain an equivalence of spectra.

\begin{introthm}[Theorems~\ref{thm:paper-local-norm} and~\ref{thm:paper-local-spectrum}]\label{thm:intro-local}
Let $F$ be a nonarchimedean local field, and let $A$ be a central simple algebra over $F$.
There is an equivalence of spectra
\[
 \nu\colon\tau_{\ge1}K(A)\xrightarrow{\ \sim\ }\tau_{\ge1}K(F).
\]
On $\pi_j$ it induces $\Nrd_j$ for $j=1,2$.
\end{introthm}


\begin{introthm}[Theorems~\ref{thm:paper-global-comparison}, \ref{thm:paper-global-split}, \ref{thm:paper-global-function-split}, \ref{thm:paper-global-fiber-bounds} and~\ref{thm:paper-global-fiber-gem}, Corollaries~\ref{cor:global-fibre-connectivity} and~\ref{cor:paper-global-fiber-prime-four}]\label{thm:intro-global}
Let $F$ be a global field and let $A$ be a central simple algebra over $F$.
If $A$ splits at every real place of $F$, all statements are integral; if $A$ is ramified at some real place, all spectra and groups are taken with $\Z[1/2]$-coefficients, and the prime $\ell=2$ is excluded.
\begin{enumerate}
\item There is a map of spectra
\[
 g\colon\tau_{\ge1}K(A)\longrightarrow\tau_{\ge1}K(F)
\]
that agrees with $\Nrd_1$ and $\Nrd_2$ in degrees one and two.
\item Let $Z=\fib(g)$.
There is an equivalence
\[
 \tau_{\ge2}K(A)\simeq\tau_{\ge2}K(F)\oplus Z
\]
whose first component is $\tau_{\ge2}g$, and $Z\in\mathrm{Sp}_{\ge4}$.
Thus, for $j\ge2$, the map $g_j=\pi_jg$ is split surjective with kernel $\pi_jZ$, and $g_2$ and $g_3$ are isomorphisms.
\item For a prime $\ell\mid\ind A$ with $\ell\ne\operatorname{char}F$, put $e_\ell=\ell^{v_\ell(\ind A)}$, and let $B_\ell(A)$ be the fiber of the localization $K(A)_{(\ell)}\to L_{E(1),\ell}K(A)$.
Then
\[
 Z\simeq\prod_{\ell}\tau_{\ge2}B_\ell(A),\qquad
 e_\ell\,\pi_jB_\ell(A)=0\ \ (j\ge1),\qquad
 B_\ell(A)\in\mathrm{Sp}_{\le e_\ell-3},
\]
and $\tau_{\ge2}B_\ell(A)=0$ if $e_\ell=\ell$.
\item
If $F$ is a global function field, then for $j\ge2$
\[
 \pi_jZ\cong K_j(A)\Big/\sum_{L}\theta_LK_j(L),
\]
where $L$ runs over the finite separable extensions of $F$ that split $A$ and $\theta_L\colon K(L)\to K(A)$ is the splitting transfer.
\item If $F$ is a number field, there is a noncanonical equivalence
\[
 Z\simeq\bigoplus_{\ell\mid\ind A}\ \bigoplus_{j=4}^{e_\ell-3}\Sigma^jH\bigl(\pi_jZ_{(\ell)}\bigr),
\]
where each $\pi_jZ_{(\ell)}$ is a finite $\ell$-group, namely the kernel of $(g_j)_{(\ell)}\colon K_j(A)_{(\ell)}\to K_j(F)_{(\ell)}$.
\end{enumerate}
\end{introthm}

Consequently, the $\ell$-primary part of $\pi_jZ$ can be nonzero only if $4\le j\le e_\ell-3$ and $e_\ell\ne\ell$.
In particular, $g$ is an equivalence if $e_\ell\in\{\ell,4\}$ for every prime $\ell\mid\ind A$ with $\ell\ne\operatorname{char}F$.

Inverting $2$ is necessary: an explicit counterexample is given in Example~\ref{ex:paper-global-k3-real-ramification}.

In degree three the comparison is explicit, and the following descriptions are integral, also in the presence of real ramification (Theorem~\ref{thm:global-k3} and Corollary~\ref{cor:global-k3-structure}).
If $F$ is a global function field with constant field $\F_q$, then $K_3(A)\cong K_3(F)\cong\Z/(q^2-1)$.
If $F$ is a number field, then $K_3(A)\cong K_3(F)/J_F^{\mathrm{ram}}$, where $J_F^{\mathrm{ram}}$ is the part of the image of $K_3^M(F)\to K_3(F)$ supported at the real places at which $A$ is ramified.
Explicitly, let $r_2$ be the number of complex places of $F$, let $r_{\mathrm{sp}}$ be the number of real places at which $A$ splits, and put $w_2=|H^0(F,\Q/\Z(2))|$; then $K_3(A)$ has free rank $r_2$, and its torsion subgroup is $\Z/w_2$ if $r_{\mathrm{sp}}=0$ and $\Z/2w_2\oplus(\Z/2)^{r_{\mathrm{sp}}-1}$ if $r_{\mathrm{sp}}>0$.


Up to equivalence, the fiber $Z$ depends only on the cyclic subgroup of $\Br(F)$ generated by $[A]$.


\begin{introthm}[Proposition~\ref{prop:paper-brauer-power-fiber}]\label{thm:intro-fiber}
Let $F$ be a global field and $A$ a central simple algebra over $F$, use the coefficient convention of Theorem~\ref{thm:intro-global}, and write $Z_A$ for the fiber of the comparison $g$ of Theorem~\ref{thm:intro-global} for $A$.
If $a\ge1$ and $\gcd(a,\per A)=1$, then there is a noncanonical equivalence $Z_A\simeq Z_{A'}$ for every central simple algebra $A'$ over $F$ with $[A']=a[A]$ in $\Br(F)$.
\end{introthm}

\begin{introthm}[Merkurjev--Suslin conjecture for higher local fields, Theorems~\ref{thm:second-two-surj}, \ref{thm:second-two-vanishing} and~\ref{thm:second-three-coker}, Corollaries~\ref{cor:review-two-local-sk1-zero} and~\ref{cor:second-three-charzero}]\label{thm:intro-higher}
Let $F$ be a two- or three-dimensional local field, let $A$ be a central simple algebra over $F$, put $\alpha=[A]\in\Br(F)$, and let $N_2(\alpha)\subseteq K_2(F)$ be the subgroup generated by the norms $N_{E/F}K_2(E)$ from all finite extensions $E/F$ that split $\alpha$.
\begin{enumerate}
\item If $F$ is two-dimensional, then $\Nrd_1$ is injective, $N_2(\alpha)=K_2(F)$, and $\Nrd_2\colon K_2(A)\to K_2(F)$ is an isomorphism.
\item If $F$ is three-dimensional, then
\[
 N_2(\alpha)=\Nrd_2K_2(A)=\ker\bigl(\alpha\cup-\colon K_2(F)\to H^4(F,\Q/\Z(3))\bigr),
\]
and $\inv_F(\alpha\cup-)$ induces an isomorphism $\coker(\Nrd_2)\cong\frac1{\per\alpha}\Z/\Z$.
If moreover the first residue field of $F$ has characteristic zero, then $\Nrd_2$ is injective.
\end{enumerate}
\end{introthm}

Together with Theorem~\ref{thm:intro-K2}, Theorem~\ref{thm:intro-higher} shows that $\SK_2(A)=\ker(\Nrd_2)$ vanishes over every local field of dimension at most two (in dimension zero, $A$ is split by Wedderburn's theorem).
In dimension four and higher the vanishing fails (Remark~\ref{rem:paper-sk2-dimension-four}).
This parallels the behavior of $\SK_1(A)=\ker(\Nrd_1)$, which vanishes over local fields of dimension at most two (\cite{NakayamaMatsushima1943,HazratWadsworth} and Theorem~\ref{thm:intro-higher}(1)) but not in dimension three \cite{Platonov1976}.
We conjecture that $\SK_2(A)=0$ over every three-dimensional local field (Conjecture~\ref{conj:paper-three-local-sk2}) and two- and three-dimensional global fields (Conjecture~\ref{conj:paper-global-sk2}).

We apply these results to noncommutative class field theory.

\begin{introthm}[Theorems~\ref{thm:paper-lca-comparison}, \ref{thm:paper-lca-fiber-splitting} and~\ref{thm:paper-lca-artin}, Corollary~\ref{cor:paper-lca-gem}]\label{thm:intro-LCA}
Let $F$ be a number field, let $A$ be a central simple algebra over $F$, and use the coefficient convention of Theorem~\ref{thm:intro-global}.
\begin{enumerate}
\item The comparisons of Theorems~\ref{thm:intro-local} and~\ref{thm:intro-global} induce a map $c^{\mathrm{lc}}_A\colon\tau_{\ge1}K(\operatorname{LCA}_A)\to\tau_{\ge1}K(\operatorname{LCA}_F)$, compatible with the localization sequences of \cite{BHRAdeles}, such that $\fib(c^{\mathrm{lc}}_A)\simeq\Sigma Z$.
There is a noncanonical equivalence
\[
 \tau_{\ge2}K(\operatorname{LCA}_A)\simeq\tau_{\ge2}K(\operatorname{LCA}_F)\oplus\Sigma Z.
\]
In particular, $c^{\mathrm{lc}}_A$ is an isomorphism in degrees $1\le n\le4$.
\item For every prime $p$, odd if $A$ is ramified at a real place, the comparison $c^{\mathrm{lc}}_A$ is compatible with Clausen's pro-$p$ Artin maps \cite{ClausenArtin}, and the Artin map of $A$ vanishes on $\fib(c^{\mathrm{lc}}_A)$.
On id\`eles of $A$, the resulting Artin map sends $(a_v)_v$ to $\operatorname{rec}_{F,p}\bigl((\Nrd_{A_v/F_v}(a_v))_v\bigr)$, where $\operatorname{rec}_{F,p}$ is the global Artin map followed by the projection to the maximal pro-$p$ quotient.
\end{enumerate}
\end{introthm}

Note that the Merkurjev--Suslin conjecture implies the isomorphism $K_2(\operatorname{LCA}_A)\cong K_2(\operatorname{LCA}_F)$ \cite[Theorem 10.5]{BHRAdeles} (Theorem~\ref{thm:paper-lca-k2}).


\subsection*{Organization}

Section~\ref{sec:local} treats nonarchimedean local fields, comparing $K(D)$ with its \'etale hypersheafification by the methods of Clausen and Mathew \cite{ClausenMathew}, with input from \cite{SuslinYufryakov}, \cite{HesselholtLarsenLindenstrauss2020}, and a Frobenius comparison of twisted cycle complexes (Appendix~\ref{app:second-frobenius}).
Section~\ref{sec:splitting-transfers} proves Theorem~\ref{thm:intro-K2} over global fields from the norm theorem of Bak and Rehmann \cite{BakRehmann1984}, cyclic Hilbert~90 for $K_2$ \cite[Theorem 14.1]{MerkurjevSuslin1982}, and a local--global principle for $K_2(D)$ modulo splitting transfers, proved with Picard norms on curves.
Section~\ref{sec:global} constructs the global comparison with the twisted motivic spectral sequence of Kahn and Levine \cite{KahnLevine2010}, arithmetic squares, transfers from homotopy orbits, Selmer $K$-theory \cite{ClausenMathew}, and Quillen's decomposition of the $K$-theory of Severi--Brauer varieties.
Sections~\ref{sec:two-local} and~\ref{sec:three-local} use Kato's higher local class field theory \cite{KatoExistence,KatoLocalII,Kurihara2000} and, for wild ramification, relative topological cyclic homology of graded covers (Appendix~\ref{app:second-tc}).
Section~\ref{sec:lca} treats locally compact modules with the localization sequences of \cite{BHRAdeles} and Clausen's Selmer $K$-homology \cite{ClausenArtin}.
Appendix~\ref{app:second-frobenius} treats fields $F$ of characteristic $p$ with $[F:F^p]=p$: it compares twisted and ordinary cycle complexes when the period is a power of $p$ and shows that $K_i(D;\Z/p^a)=0$ for $i\ge2$, which is the input at the characteristic prime in Sections~\ref{sec:local}--\ref{sec:global}.
Appendix~\ref{app:second-tc} compares relative $K$-theory of maximal orders with relative $\mathrm{TC}$ at finite truncations and represents the resulting classes by integral classes, as needed for wild ramification in Section~\ref{sec:two-local} and for deep principal units in Section~\ref{sec:three-local}.
Appendix~\ref{app:review-cycle-transfer} collects facts on zero-cycles with coefficients.
